% 第6章 6.2.2：拉普拉斯读作“本点值减邻居平均，再乘度”。单位权六点路径，f=(0,1,2,4,3,3)。
% 邻居平均 (1,1,2.5,2.5,3.5,3)，差 f-平均 (-1,0,-0.5,1.5,-0.5,0)，度 (1,2,2,2,2,1)，
% Lf=(-1,0,-1,3,-1,0)，能量 f^T L f=7（各边差 1,1,2,1,0 的平方和）。数值由 c6-path-eight-modes/calc.py 复算。
\begin{tikzpicture}[x=1mm,y=1mm,>=Stealth,font=\figurefont\footnotesize,text=FigureInk,
  vertex/.style={circle,draw=FigureStroke,fill=FigureBlueFill,line width=.8pt,minimum size=5.6mm,inner sep=0pt},
  edge/.style={draw=FigureStroke,line width=.8pt},
  val/.style={circle,fill=FigureBlue,draw=FigureBlue,inner sep=0pt,minimum size=1.8mm},
  avg/.style={circle,draw=FigureYellow,fill=white,line width=.8pt,inner sep=0pt,minimum size=3.2mm},
  dev/.style={->,draw=FigureStroke,line width=.6pt,shorten >=1pt,shorten <=1.8pt},
  note/.style={font=\FigureNoteFont\fontsize{7.5}{9}\selectfont,text=FigureMuted}]
\def\ys{8}
% 坐标轴与网格
\foreach \t in {0,1,2,3,4}{\draw[FigureGrid,line width=.4pt] (4,\t*\ys+22)--(82,\t*\ys+22);
  \node[anchor=east,font=\figurefont\scriptsize] at (3,\t*\ys+22) {$\t$};}
% 数值折线
\draw[FigureBlue,line width=.6pt] \foreach \v [count=\i] in {0,1,2,4,3,3}{\ifnum\i=1 (\i*14-5,\v*\ys+22)\else --(\i*14-5,\v*\ys+22)\fi};
% 邻居平均与偏差箭头
\foreach \v/\a [count=\i] in {0/1,1/1,2/2.5,4/2.5,3/3.5,3/3}{
  \pgfmathsetmacro{\yv}{\v*\ys+22}\pgfmathsetmacro{\ya}{\a*\ys+22}
  \ifdim\v pt=\a pt\else\draw[dev] (\i*14-5,\ya)--(\i*14-5,\yv);\fi
  \node[avg] at (\i*14-5,\ya) {};
  \node[val] at (\i*14-5,\yv) {};}
% 路径与 Lf
\foreach \i in {1,...,6}{\node[vertex] (v\i) at (\i*14-5,10) {$\i$};}
\foreach \i [evaluate=\i as \j using int(\i+1)] in {1,...,5}{\draw[edge] (v\i)--(v\j);}
\foreach \y [count=\i] in {1,2,2,2,2,1}{\node at (\i*14-5,3) {$\y$};}
\foreach \y [count=\i] in {-1,0,-1,3,-1,0}{\node at (\i*14-5,-3) {$\y$};}
\node[anchor=east] at (3,10) {顶点};
\node[anchor=east] at (3,3) {度$d_i$};
\node[anchor=east] at (3,-3) {$(\bm L\bm f)_i$};
% 图例
\node[val] at (6,59) {}; \node[anchor=west,note] at (8,59) {本点值$f_i$};
\node[avg] at (30,59) {}; \node[anchor=west,note] at (32,59) {邻居平均$\bar f_i$};
\draw[dev] (58,57)--(58,61); \node[anchor=west,note] at (60,59) {差$f_i-\bar f_i$};
\end{tikzpicture}
